% werte.tex — die Werteschnittstelle. Nur diese vier Makros, sonst nichts.
%
% Sie stand wortgleich in mechanik.tex und in mappe/mechanik-mappe.tex, mit
% dem Kommentar „Gleiche Schnittstelle wie mechanik.tex, damit
% helden/<name>.tex unverändert durch beide Dokumente läuft". Genau das ist
% der Grund, warum sie nicht zweimal stehen darf: die beiden Fassungen waren
% schon auseinandergelaufen — \wertoder gab es nur in der Mappe, und \wertdaP
% war einmal einzeilig und einmal umbrochen geschrieben. Solange sie sich nur
% im Zeilenumbruch unterscheiden, fällt es nicht auf; beim nächsten Zusatz
% schon.
%
% Eingebunden mit % werte.tex — die Werteschnittstelle. Nur diese vier Makros, sonst nichts.
%
% Sie stand wortgleich in mechanik.tex und in mappe/mechanik-mappe.tex, mit
% dem Kommentar „Gleiche Schnittstelle wie mechanik.tex, damit
% helden/<name>.tex unverändert durch beide Dokumente läuft". Genau das ist
% der Grund, warum sie nicht zweimal stehen darf: die beiden Fassungen waren
% schon auseinandergelaufen — \wertoder gab es nur in der Mappe, und \wertdaP
% war einmal einzeilig und einmal umbrochen geschrieben. Solange sie sich nur
% im Zeilenumbruch unterscheiden, fällt es nicht auf; beim nächsten Zusatz
% schon.
%
% Eingebunden mit % werte.tex — die Werteschnittstelle. Nur diese vier Makros, sonst nichts.
%
% Sie stand wortgleich in mechanik.tex und in mappe/mechanik-mappe.tex, mit
% dem Kommentar „Gleiche Schnittstelle wie mechanik.tex, damit
% helden/<name>.tex unverändert durch beide Dokumente läuft". Genau das ist
% der Grund, warum sie nicht zweimal stehen darf: die beiden Fassungen waren
% schon auseinandergelaufen — \wertoder gab es nur in der Mappe, und \wertdaP
% war einmal einzeilig und einmal umbrochen geschrieben. Solange sie sich nur
% im Zeilenumbruch unterscheiden, fällt es nicht auf; beim nächsten Zusatz
% schon.
%
% Eingebunden mit % werte.tex — die Werteschnittstelle. Nur diese vier Makros, sonst nichts.
%
% Sie stand wortgleich in mechanik.tex und in mappe/mechanik-mappe.tex, mit
% dem Kommentar „Gleiche Schnittstelle wie mechanik.tex, damit
% helden/<name>.tex unverändert durch beide Dokumente läuft". Genau das ist
% der Grund, warum sie nicht zweimal stehen darf: die beiden Fassungen waren
% schon auseinandergelaufen — \wertoder gab es nur in der Mappe, und \wertdaP
% war einmal einzeilig und einmal umbrochen geschrieben. Solange sie sich nur
% im Zeilenumbruch unterscheiden, fällt es nicht auf; beim nächsten Zusatz
% schon.
%
% Eingebunden mit \input{werte} — das Arbeitsverzeichnis ist bogen/, dafür
% sorgen die Bauskripte.
%
% \makeatletter hier selbst, nicht beim Aufrufer: \@firstoftwo braucht es,
% und eine Datei, die nur unter einer bestimmten Katcode-Lage funktioniert,
% ist eine Falle für den, der sie später woanders einbindet.
%
% Am Ende steht deshalb NICHT \makeatother, sondern die Lage des Aufrufers,
% wie sie beim Eintritt war. Der Unterschied ist einmal aufgetreten und
% kostet Suchzeit: mechanik.tex steht selbst unter \makeatletter, nach einem
% \makeatother wäre dort das @ kein Buchstabe mehr, und der Lauf endete mit
%
%   ! Undefined control sequence.
%   \feld @zeile->\endgroup \g @addto@macro\feldliste {...}
%
% — mit einem Zeilenverweis auf felder/df.tex. Die Meldung zeigt also auf
% die Feldtabelle, während der Fehler hier steht.

\expandafter\edef\csname werte@katcode\endcsname{%
  \catcode`\noexpand\@=\the\catcode`\@\relax}
\makeatletter

% \wert{feldname}{Wert} kommt aus helden/<name>.tex. Ein fehlender Wert ist
% kein Fehler, sondern ein leeres Feld — der Lauf darf daran nicht scheitern.
\newcommand{\wert}[2]{\expandafter\gdef\csname wert@#1\endcsname{#2}}

% Der Wert, oder nichts.
\newcommand{\wertvon}[1]{\ifcsname wert@#1\endcsname\csname wert@#1\endcsname\fi}

% \wertdaP{name}{dann}{sonst} — das P steht für Prädikat. Verzweigt, ohne den
% Wert selbst zu setzen.
\newcommand{\wertdaP}[1]{%
  \ifcsname wert@#1\endcsname\expandafter\@firstoftwo\else\expandafter\@secondoftwo\fi}

% Der erste gesetzte Wert aus zwei Namen, sonst leer. Damit greift die Mappe
% auf die Felder des Heldenbogens zurück: mappe_name, sonst s1_name.
\newcommand{\wertoder}[2]{\wertdaP{#1}{\wertvon{#1}}{\wertvon{#2}}}

\werte@katcode
 — das Arbeitsverzeichnis ist bogen/, dafür
% sorgen die Bauskripte.
%
% \makeatletter hier selbst, nicht beim Aufrufer: \@firstoftwo braucht es,
% und eine Datei, die nur unter einer bestimmten Katcode-Lage funktioniert,
% ist eine Falle für den, der sie später woanders einbindet.
%
% Am Ende steht deshalb NICHT \makeatother, sondern die Lage des Aufrufers,
% wie sie beim Eintritt war. Der Unterschied ist einmal aufgetreten und
% kostet Suchzeit: mechanik.tex steht selbst unter \makeatletter, nach einem
% \makeatother wäre dort das @ kein Buchstabe mehr, und der Lauf endete mit
%
%   ! Undefined control sequence.
%   \feld @zeile->\endgroup \g @addto@macro\feldliste {...}
%
% — mit einem Zeilenverweis auf felder/df.tex. Die Meldung zeigt also auf
% die Feldtabelle, während der Fehler hier steht.

\expandafter\edef\csname werte@katcode\endcsname{%
  \catcode`\noexpand\@=\the\catcode`\@\relax}
\makeatletter

% \wert{feldname}{Wert} kommt aus helden/<name>.tex. Ein fehlender Wert ist
% kein Fehler, sondern ein leeres Feld — der Lauf darf daran nicht scheitern.
\newcommand{\wert}[2]{\expandafter\gdef\csname wert@#1\endcsname{#2}}

% Der Wert, oder nichts.
\newcommand{\wertvon}[1]{\ifcsname wert@#1\endcsname\csname wert@#1\endcsname\fi}

% \wertdaP{name}{dann}{sonst} — das P steht für Prädikat. Verzweigt, ohne den
% Wert selbst zu setzen.
\newcommand{\wertdaP}[1]{%
  \ifcsname wert@#1\endcsname\expandafter\@firstoftwo\else\expandafter\@secondoftwo\fi}

% Der erste gesetzte Wert aus zwei Namen, sonst leer. Damit greift die Mappe
% auf die Felder des Heldenbogens zurück: mappe_name, sonst s1_name.
\newcommand{\wertoder}[2]{\wertdaP{#1}{\wertvon{#1}}{\wertvon{#2}}}

\werte@katcode
 — das Arbeitsverzeichnis ist bogen/, dafür
% sorgen die Bauskripte.
%
% \makeatletter hier selbst, nicht beim Aufrufer: \@firstoftwo braucht es,
% und eine Datei, die nur unter einer bestimmten Katcode-Lage funktioniert,
% ist eine Falle für den, der sie später woanders einbindet.
%
% Am Ende steht deshalb NICHT \makeatother, sondern die Lage des Aufrufers,
% wie sie beim Eintritt war. Der Unterschied ist einmal aufgetreten und
% kostet Suchzeit: mechanik.tex steht selbst unter \makeatletter, nach einem
% \makeatother wäre dort das @ kein Buchstabe mehr, und der Lauf endete mit
%
%   ! Undefined control sequence.
%   \feld @zeile->\endgroup \g @addto@macro\feldliste {...}
%
% — mit einem Zeilenverweis auf felder/df.tex. Die Meldung zeigt also auf
% die Feldtabelle, während der Fehler hier steht.

\expandafter\edef\csname werte@katcode\endcsname{%
  \catcode`\noexpand\@=\the\catcode`\@\relax}
\makeatletter

% \wert{feldname}{Wert} kommt aus helden/<name>.tex. Ein fehlender Wert ist
% kein Fehler, sondern ein leeres Feld — der Lauf darf daran nicht scheitern.
\newcommand{\wert}[2]{\expandafter\gdef\csname wert@#1\endcsname{#2}}

% Der Wert, oder nichts.
\newcommand{\wertvon}[1]{\ifcsname wert@#1\endcsname\csname wert@#1\endcsname\fi}

% \wertdaP{name}{dann}{sonst} — das P steht für Prädikat. Verzweigt, ohne den
% Wert selbst zu setzen.
\newcommand{\wertdaP}[1]{%
  \ifcsname wert@#1\endcsname\expandafter\@firstoftwo\else\expandafter\@secondoftwo\fi}

% Der erste gesetzte Wert aus zwei Namen, sonst leer. Damit greift die Mappe
% auf die Felder des Heldenbogens zurück: mappe_name, sonst s1_name.
\newcommand{\wertoder}[2]{\wertdaP{#1}{\wertvon{#1}}{\wertvon{#2}}}

\werte@katcode
 — das Arbeitsverzeichnis ist bogen/, dafür
% sorgen die Bauskripte.
%
% \makeatletter hier selbst, nicht beim Aufrufer: \@firstoftwo braucht es,
% und eine Datei, die nur unter einer bestimmten Katcode-Lage funktioniert,
% ist eine Falle für den, der sie später woanders einbindet.
%
% Am Ende steht deshalb NICHT \makeatother, sondern die Lage des Aufrufers,
% wie sie beim Eintritt war. Der Unterschied ist einmal aufgetreten und
% kostet Suchzeit: mechanik.tex steht selbst unter \makeatletter, nach einem
% \makeatother wäre dort das @ kein Buchstabe mehr, und der Lauf endete mit
%
%   ! Undefined control sequence.
%   \feld @zeile->\endgroup \g @addto@macro\feldliste {...}
%
% — mit einem Zeilenverweis auf felder/df.tex. Die Meldung zeigt also auf
% die Feldtabelle, während der Fehler hier steht.

\expandafter\edef\csname werte@katcode\endcsname{%
  \catcode`\noexpand\@=\the\catcode`\@\relax}
\makeatletter

% \wert{feldname}{Wert} kommt aus helden/<name>.tex. Ein fehlender Wert ist
% kein Fehler, sondern ein leeres Feld — der Lauf darf daran nicht scheitern.
\newcommand{\wert}[2]{\expandafter\gdef\csname wert@#1\endcsname{#2}}

% Der Wert, oder nichts.
\newcommand{\wertvon}[1]{\ifcsname wert@#1\endcsname\csname wert@#1\endcsname\fi}

% \wertdaP{name}{dann}{sonst} — das P steht für Prädikat. Verzweigt, ohne den
% Wert selbst zu setzen.
\newcommand{\wertdaP}[1]{%
  \ifcsname wert@#1\endcsname\expandafter\@firstoftwo\else\expandafter\@secondoftwo\fi}

% Der erste gesetzte Wert aus zwei Namen, sonst leer. Damit greift die Mappe
% auf die Felder des Heldenbogens zurück: mappe_name, sonst s1_name.
\newcommand{\wertoder}[2]{\wertdaP{#1}{\wertvon{#1}}{\wertvon{#2}}}

\werte@katcode
